\documentclass[../../main.tex]{subfiles}

\begin{document}

\chapter{Basic operation}

\section{Screens and settings}

%===========================================================
%===========================================================
%               Main Screen
%===========================================================
%===========================================================

\subsection{Main screen}

\begin{figure}[H]
    \centering
    \includegraphics[width=1\textwidth]{GUISuruga/01_main_screen.png}
    \caption{Main screen of the Suruga positioning system. The interface is divided into four main sections: (a) Axis Status, which displays the current position and status of each axis; (b) Optical Power, which displays the optical power measured by the connected power meter; (c) Analog Input, which displays the voltage measured at each available analog input channel; and (d) Function Access, which provides direct access to the main functionalities and configuration options of the Suruga platform.}
    \label{fig:suruga_main_screen}
\end{figure}

This screen is divided into three main sections. The first section, Axis Status, allows the user to monitor the current position of each axis, the CW and CCW limit switch status, the origin sensor status, the alarm outputs, and the availability of each servo. As listed below:

\begin{itemize}
    \item \textbf{Current Position:} Displays the current coordinates of each axis. Linear axes are displayed in micrometers ($\mu$m), while rotational axes are displayed in degrees ($^\circ$). Axes with the number 1 correspond to the left arm of the Suruga positioning system, while axes with the number 2 correspond to the right arm.

    \item \textbf{CW:} Displays the status of the clockwise (CW) limit sensor. Red indicates that the limit has been reached, while gray indicates that the axis is within the safe operating range.

    \item \textbf{CCW:} Displays the status of the counterclockwise (CCW) limit sensor. Red indicates that the limit has been reached, while gray indicates that the axis is within the safe operating range.

    \item \textbf{ORG:} Displays the status of the origin sensor. For this particular stage, no origin sensor is available, so the indicator will remain continuously displayed in red or gray.

    \item \textbf{ALM:} Indicates the alarm status of the axis. Red is displayed when an alarm occurs, such as an overload or a similar fault condition. Gray indicates normal operation.

    \item \textbf{RDY:} Displays the operation status of each axis. Green indicates that the servo is enabled and ready for operation, while gray indicates that the servo is turned off.
\end{itemize}

The second section, Optical Power, displays the current reading of the power meter, if one is available. Finally, the Analog Input section allows the user to monitor the values of the different analog inputs. On the left side of the screen, direct access to the different functionalities of the Suruga platform is provided.

%===========================================================
%===========================================================
%               Jog screen
%===========================================================
%===========================================================

\subsection{Jog screen}

By pressing the Jog button in the quick access section, indicated in part (d) of the main screen in Fig.~\ref{fig:suruga_main_screen}(d), the user is directed to the Jog screen shown in Fig.~\ref{fig:02_jog_screen}.

This screen allows the user to perform movements along the different axes of the stage and to adjust the rotation angles of the corresponding axes. It is important to note that, once the Suruga Seiki stage is powered on, no power is supplied to the motors by default. Therefore, the Servo button on this screen must be pressed to enable each axis of the stage before any movement can be performed.

\begin{figure}[H]
    \centering
    \includegraphics[width=1\textwidth]{GUISuruga/02_00_jog_screen.png}
    \caption{Jog screen of the Suruga positioning system. The interface is divided into three main sections: (a) Measurement Monitoring, which displays the optical power and analog input measurements; (b) Movement Control, which provides the servo, error reset, movement mode, velocity, and stop controls; and (c) Axis Control, which provides independent movement and homing controls for each axis of both arms, together with the corresponding position, velocity, limit switch, and alarm status information.}
    \label{fig:02_jog_screen}
\end{figure}

The jog screen can be divided into three sections. Section (a) Measurement Monitoring, displays the optical power measured by the power meter, as well as the voltage measured through the analog inputs. Section (b)  Movement Control, provides access to five control buttons and selectors:

\begin{itemize}
    \item \textbf{Servo:} Enables or disables the servos of the stage.

    \item \textbf{Error Reset:} Resets the servo alarm in case an alarm
    occurs.

    \item \textbf{Mode:} Allows the user to select between the three
    available movement modes: Jog, Inching, and
    Target. In Jog mode, the movement continues while
    the movement button is held down. In Inching mode, pressing
    a movement button moves the selected axis by the distance specified
    in the Inching Distance parameter. In Target mode,
    the selected axis moves to the target position specified in Section
    (c).

    \item \textbf{Velocity} -- Allows the user to select between the three
    preset speed configurations: High, Mid, and Low.

    \item \textbf{Stop} -- Stops the movement of the selected axis.
\end{itemize}


Finally, Section (c) Axis Control, is divided into twelve subsections, six for each arm. Each subsection corresponds to one of the six degrees of freedom available for the respective arm: translational movements along the $X$, $Y$, and $Z$ axes, and angular movements around the $X$, $Y$, and $Z$ axes ($T_x$, $T_y$, and $T_z$, respectively). Each of the subsections in Section (c) provides a set of movement control buttons, as well as fields for defining the target position and the inching distance. The user can also define the movement velocity, while the current position of the corresponding axis is continuously displayed. In addition, each subsection includes status indicators for the CW, CCW, and ORG limit switches, as well as an ALM indicator.

\begin{figure}[H]
    \centering
    \includegraphics[width=1\textwidth]{GUISuruga/02_01_move_buttons.png}
    \caption{Movement controls and status indicators for an individual stage axis, including target position, inching distance, velocity, current position, limit-switch status, alarm status, and homing control.}
    \label{fig:02_01_move_buttons}
\end{figure}

The following figures illustrate the movement directions available for the
three translational degrees of freedom of the stage. The translational
movements along the $X$, $Y$, and $Z$ axes are presented, showing the
directions associated with the \texttt{+} and \texttt{-} buttons for each
arm.

\begin{figure}[H]
    \centering
    \includegraphics[width=1\textwidth]{GUISuruga/02_02_JogX.png}
    \caption{Movement directions along the $X$ axis for arm 1 (a) and arm 2 (b) associated with the \texttt{+} and \texttt{-} buttons.}
    \label{fig:02_02_JogX}
\end{figure}

\begin{figure}[H]
    \centering
    \includegraphics[width=1\textwidth]{GUISuruga/02_03_JogY.png}
    \caption{Movement directions along the $Y$ axis for arm 1 (a) and arm 2 (b) associated with the \texttt{+} and \texttt{-} buttons.}
    \label{fig:02_03_JogY}
\end{figure}

\begin{figure}[H]
    \centering
    \includegraphics[width=1\textwidth]{GUISuruga/02_04_JogZ.png}
    \caption{Movement directions along the $Z$ axis for arm 1 (a) and arm 2 (b) associated with the \texttt{+} and \texttt{-} buttons.}
    \label{fig:02_04_JogZ}
\end{figure}

The corresponding angular movements are presented next. These figures
illustrate the rotation directions associated with the \texttt{+} and
\texttt{-} buttons for rotations around the $X$, $Y$, and $Z$ axes
($T_x$, $T_y$, and $T_z$, respectively), for both arms of the stage.


\begin{figure}[H]
    \centering
    \includegraphics[width=1\textwidth]{GUISuruga/JogTx.png}
    \caption{Rotation directions around the $X$ axis produced by the \texttt{+} and \texttt{-} buttons for arm 1 (a) and arm 2 (b).}
    %\label{fig:frontDA1000}
\end{figure}


\begin{figure}[H]
    \centering
    \includegraphics[width=1\textwidth]{GUISuruga/JogTy.png}
    \caption{Pendiente.}
    \caption{Rotation directions around the $Y$ axis produced by the \texttt{+} and \texttt{-} buttons for arm 1 (a) and arm 2 (b).}
\end{figure}

\begin{figure}[H]
    \centering
    \includegraphics[width=1\textwidth]{GUISuruga/JogTz.png}
    \caption{Rotation directions around the $Z$ axis produced by the \texttt{+} and \texttt{-} buttons for arm 1 (a) and arm 2 (b).}
    %\label{fig:frontDA1000}
\end{figure}



%===========================================================
%===========================================================
%               Profile screen
%===========================================================
%===========================================================

\subsection{Profile screen}

The SurugaSeiki.Motion.Profile class provides a high-level interface for configuring, executing, and monitoring profiling operations on the Suruga Seiki platform. To use this functionality, an instance of the Profile class is created, providing access to a set of methods for configuring the acquisition parameters, starting and stopping the profiling process, monitoring its status, and requesting the acquired data.

The class also contains several nested types that complement its functionality. ProfileParameter is used to define the parameters required to configure a profiling operation, while ProfileData represents the data acquired during the operation. In addition, the Status and ErrorCode enumerations are used to identify, respectively, the current status of the profiling process and the error conditions that may occur during its execution.

For a detailed description of the functionalities, methods, and data structures of the Profile class, refer to the class definition provided in Appendix~\ref{sec:profile_class}. However in this section we will talk about the Profile screen.

\begin{figure}[H]
    \centering
    \includegraphics[width=1\textwidth]{GUISuruga/03_profile_screen.png}
    \caption{Pendiente.}
    %\label{fig:frontDA1000}
\end{figure}



%===========================================================
%===========================================================
%               Optical alignment screen
%===========================================================
%===========================================================

\subsection{Optical alignment screen}

\begin{figure}[H]
    \centering
    \includegraphics[width=1\textwidth]{GUISuruga/04_optical_alignment_screen.png}
    \caption{Pendiente.}
    %\label{fig:frontDA1000}
\end{figure}

El usuario puede llevar a cabo los siguientes métodos de alineamiento:

\begin{itemize}
    \item \textbf{Alineamiento Single (Single Alignment):} Realiza el alineamiento utilizando un único eje. 
    \item \textbf{Alineamiento Flat (Flat Alignment):} Realiza el alineamiento en un plano utilizando dos ejes. 
    \item \textbf{Alineamiento Focus (Focus Alignment):} Además del alineamiento Flat, realiza el ajuste sobre el eje óptico empleando tres ejes. 
    \item \textbf{Alineamiento por rotación de dos canales (2ch Rotation Alignment):} Realiza un alineamiento Flat utilizando dos señales ópticas y corrige el ángulo alrededor del eje óptico en función de la posición y de la distancia configurada. 
    \item \textbf{Alineamiento MMF (MMF Rotation Alignment):} Posiciona una fibra multimodo (MMF, Multi-Mode Fiber) en el lado de salida y corrige la posición y el ángulo del arreglo de fibras en el lado de entrada. 
\end{itemize}

%===========================================================
%===========================================================

\subsubsection{Alineamiento Flat (Flat Alignment):}

\begin{figure}[H]
    \centering
    \includegraphics[width=1\textwidth]{Alignment/FlatAlignment.png}
    \caption{Pendiente.}
    %\label{fig:frontDA1000}
\end{figure}

Se parte de la condición en la que la fibra y el Device Under Test (DUT) se encuentran desalineados, por lo que el sistema desconoce la ubicación del punto de máxima potencia óptica. Debido a ello, no es posible iniciar directamente un Peak Search, ya que este requiere partir de una señal detectable para localizar el máximo.

Como primer paso, el algoritmo ejecuta un Field Search, cuyo objetivo es explorar el área alrededor de la posición inicial hasta encontrar una potencia lo suficientemente grande. La condición que determina el éxito de esta etapa es:

\begin{equation}
P_{\mathrm{OPM}} > \mathrm{FieldSearchThreshold}
\label{eq:field_search_threshold}
\end{equation}

Donde FieldSearchThreshold corresponde al nivel mínimo de potencia, expresada en voltios haciendo uso de la salida analógica del OPM, que indica que el sistema ha encontrado una señal válida.

Durante esta etapa, el movimiento sigue una trayectoria en espiral cuadrada que se expande progresivamente alrededor del punto inicial. En cuanto la potencia medida supera el umbral establecido, el algoritmo abandona el Field Search e inicia el proceso de Peak Search, encargado de localizar con precisión el punto de máxima potencia.

Si la trayectoria alcanza los límites definidos por SearchRangeX o SearchRangeY sin encontrar una señal superior a FieldSearchThreshold, el proceso finaliza con el estado FieldSearchRangeOver, indicando que no fue posible localizar una señal dentro del área de búsqueda especificada.

El comportamiento del algoritmo empleado para un alineamiento tipo flat está determinado por un conjunto de parámetros que definen el área de búsqueda, la resolución del barrido y el criterio para iniciar la búsqueda fina. En la Tabla \ref{tab:parameters} se describen los parámetros más relevantes en conjunto con su equivalente en el DLL provisto por Suruga Seiki.


\begin{longtblr}[
    caption={Parámetros de configuración},
    label={tab:parameters},
]{
    width=\textwidth,
    colspec={|Q[c]|X|Q[c]|},
    hlines,
}

\SetRow{bg=gray9}
\textbf{Parámetro} &
\textbf{Descripción} &
\textbf{Tipo .NET} \\

%===========================================================
% Alignment
%===========================================================
\SetRow{bg=SkyBlue}
\SetCell[c=3]{c}
Alignment \\

mainStageNumberX &
Número del eje principal utilizado para el movimiento en X. &
UInt16 \\

mainStageNumberY &
Número del eje principal utilizado para el movimiento en Y. &
UInt16 \\

subStageNumberXY & 
Número del sub-eje utilizado para el movimiento en XY. & 
UInt16 \\

subAngleX & 
Ángulo del sub-eje para la alineación en X. & 
Double \\

subAngleY & 
Ángulo del sub-eje para la alineación en Y. & 
Double \\

%===========================================================
% Sensor
%===========================================================
\SetRow{bg=SweetSalmon}
\SetCell[c=3]{c}
Sensor \\

pmCh & Canal del medidor de potencia utilizado durante la alineación. & UInt16 \\

analogCh & Canal de entrada analógica utilizado para la medición. & UInt16 \\

pmAutoRangeUpOn & Habilita el incremento automático del rango del medidor de potencia. & Boolean \\

pmInitRangeSettingOn & Habilita el uso de un rango inicial para el medidor de potencia. & Boolean \\

pmInitRange & Valor del rango inicial del medidor de potencia. & Int16 \\

%===========================================================
% Threshold
%===========================================================
\SetRow{bg=Periwinkle}
\SetCell[c=3]{c}
Threshold \\

fieldSearchThreshold & Umbral utilizado durante la búsqueda de campo. & Double \\

peakSearchThreshold & Umbral utilizado durante la búsqueda del pico de señal. & Double \\

%===========================================================
% Search
%===========================================================
\SetRow{bg=MutedTeal}
\SetCell[c=3]{c}
Search \\

searchRangeX & Rango de búsqueda en la dirección X. & Double \\

searchRangeY & Rango de búsqueda en la dirección Y. & Double \\

fieldSearchPitchX & Paso de exploración en X durante la búsqueda de campo. & Double \\

fieldSearchPitchY & Paso de exploración en Y durante la búsqueda de campo. & Double \\

fieldSearchFirstPitchX & Primer paso de exploración en X durante la búsqueda de campo. & Double \\

fieldSearchSpeedX & Velocidad de exploración en X durante la búsqueda de campo. & Double \\

fieldSearchSpeedY & Velocidad de exploración en Y durante la búsqueda de campo. & Double \\

peakSearchSpeedX & Velocidad de búsqueda del pico en X. & Double \\

peakSearchSpeedY & Velocidad de búsqueda del pico en Y. & Double \\

smoothingRangeX & Ventana de suavizado aplicada a los datos en X. & UInt32 \\

smoothingRangeY & Ventana de suavizado aplicada a los datos en Y. & UInt32 \\

centroidThresholdX & Umbral para el cálculo del centroide en X. & Double \\

centroidThresholdY & Umbral para el cálculo del centroide en Y. & Double \\

%===========================================================
% Judge
%===========================================================
\SetRow{bg=ApricotCream}
\SetCell[c=3]{c}
Judge \\

convergentRangeX & Rango de convergencia permitido en X. & Double \\

convergentRangeY & Rango de convergencia permitido en Y. & Double \\

comparisonCount & Número de comparaciones realizadas para verificar la convergencia. & UInt16 \\

maxRepeatCount & Número máximo de repeticiones permitidas durante el proceso de alineación. & UInt16 \\

\end{longtblr}


\end{document}


